\begin{result}{E.2.25}{One integral comparison for all surface and CM geometric factors}{Proved}
Let \(\mathcal A_{\leq2,\mathrm{CM}}\) consist of all
abelian varieties over \(\Qb\) geometrically isogenous to
products of simple factors of dimension at most two and
arbitrary CM abelian varieties. Take the full integral
diagram of E.2.16, including every geometric isogeny
lattice, all products, algebraic correspondences, tensor
constructions and specified projectors, finite-model
Artin module maps, and all Artin objects.
With \(P\) and its root labels as in E.2.23, there is a
specified \(\sigma_*\in\Gamma_{F_{\mathrm{cm}}}\) such that
\[
 \sigma_*|_P=(12)(34),\qquad
 \rho_A(\sigma_*)\Lambda_A
       =\Lambda_{{}^{\sigma_*}A}
       \quad(A\in\mathcal A_{\leq2,\mathrm{CM}}).
 \tag{E.2.25a}
\]
These actual completed transports give one integral
tensor comparison on the entire stated diagram, with
Tate map \(1\). By E.2.3 it compares the two specified
real embeddings of \(\Q(\alpha_1)\). This is the plain
integral diagram: E.2.7 supplies the explicit obstruction
when labelled real Frobenius is additionally required.
\end{result}
\begin{proof}
We prove the exact finite-family openness and integral
coset intersection before performing the simultaneous
construction.

\paragraph{1. The actual algebraic groups.}
Let \(A_i/\Qb\) be a finite family of pairwise
geometrically nonisogenous, absolutely simple, non-CM
abelian varieties of dimension one or two. Write
\[
 V_i=H_1(A_i(\C),\Q),\qquad
 \Lambda_i=H_1(A_i(\C),\Z),
\]
and choose labelled algebraic polarizations \(\psi_i\).
Enlarge a finite field of definition to define all
geometric endomorphisms and the finite set of labelled
Hodge tensors needed below.

The simple-surface classification used here is explicit
in the proofs of Lombardo's Corollaries 4.4--4.5,
author p.~12: there is no absolutely simple surface of
type III in characteristic zero, and the endomorphism
algebra of an absolutely simple surface cannot be an
imaginary quadratic field. Thus a non-CM simple surface
has geometric endomorphism algebra \(\Q\), a real
quadratic field \(F\), or an indefinite quaternion
division algebra \(B\). The remaining simple-surface
case is CM and already has the comparison of E.2.6
\cite[Corollaries 4.4--4.5 and their proofs, p.~12]
{GenLombardoNonsimple2015}.

The derived and full Mumford--Tate groups, in their
actual actions on \(V_i\), are as follows:
\[
\begin{array}{c|c|c}
 \End^0(A_i)&H_i&M_i\\ \hline
 \Q&\operatorname{Sp}(V_i,\psi_i)
       &\operatorname{GSp}(V_i,\psi_i)\\
 F&\operatorname{Res}_{F/\Q}\operatorname{SL}_{F}(V_i)
       &\{g\in\GL_F(V_i):\det_F(g)\in\Q^\times\}\\
 B&\operatorname{SL}_1(B^{\mathrm{op}})
       &(B^{\mathrm{op}})^\times .
\end{array}
\tag{E.2.25b}
\]
For an elliptic non-CM factor the first row is
\(\operatorname{SL}_2\subset\GL_2\).

Here is the verification of the actions and multipliers
in (E.2.25b). In the real multiplication case \(V_i\)
has dimension two over \(F\), and Rosati fixes \(F\)
pointwise. Indeed the nontrivial involution of a real
quadratic field sends a nonzero trace-zero element
\(a\) to \(-a\), contradicting positivity of
\(\operatorname{Tr}(aa^*)\). The trace pairing defines
a unique \(F\)-valued form \(\Psi_i\) by
\[
 \operatorname{Tr}_{F/\Q}\bigl(a\Psi_i(v,w)\bigr)
                  =\psi_i(av,w)\qquad(a\in F).
\]
Self-adjointness of \(F\) gives \(F\)-bilinearity,
alternation and nondegeneracy. In dimension two its
similitude is its determinant. This proves the displayed
rational-multiplier condition.

In the quaternionic case \(V_i\) is a free left
\(B\)-module of rank one; its centralizer is right
multiplication by \(B^{\mathrm{op}}\). Lombardo's
Lemmas 5.3--5.4 explicitly identify these right Galois
and left endomorphism actions at a sufficiently large
split prime, and Lemma 5.6 identifies the right
determinant with the cyclotomic multiplier
\cite[Lemmas 5.3--5.6, pp.~19--20]
{GenLombardoSurfaces2015}. Faithful extension
\(\Q\to\Q_\ell\) then proves over \(\Q\) that the
multiplier of right multiplication by \(b\) is
\(\operatorname{Nrd}(b)\).

To verify that the groups are the full groups in
(E.2.25b), use the Mumford--Tate equality for each
factor from Lombardo's Corollary 4.5, and the respective
full large-prime images from Theorem 1.3, Corollary 1.5
and Theorem 1.7 of his surface paper
\cite[Theorem 1.3, p.~2; Corollary 1.5 and
Theorem 1.7, p.~3]{GenLombardoSurfaces2015}.
The Mumford--Tate group lies in the displayed
endomorphism/polarization stabilizer because these
are labelled Hodge tensors. At one sufficiently large
prime this inclusion is an equality by those two
source statements. Faithful scalar extension gives
the equality over \(\Q\).

For the product \(A=\prod_i A_i\), Lombardo's
Corollary 4.4 gives, in the actual block representation,
\[
 H(A)=H:=\prod_i H_i,\qquad
 M:=\operatorname{MT}(A)=\Gm H,
 \qquad M^{\mathrm{der}}=H.
\]
Its factor projections are the given maps on homology;
the cited theorem proves that their inclusion into
the product is onto. The passage from the stated
\(\ell\)-adic equality to the Hodge equality uses
the same corollary's Mumford--Tate assertion
\cite[Corollary 4.4, p.~12]{GenLombardoNonsimple2015}.

Each \(H_i\) is simply connected and semisimple and
contains \(-I\) in this representation. The map
\((t,h)\mapsto th\) identifies \(M\) with
\((\Gm\times H)/\mu_2\), where the kernel is the
diagonal pair \((-1,-I)\). Its common multiplier
\(\nu(th)=t^2\) therefore gives the exact sequence
of algebraic groups
\[
 1\longrightarrow H\longrightarrow M
       \xrightarrow{\nu}\Gm\longrightarrow1.
 \tag{E.2.25c}
\]
In particular its multiplier-one kernel is precisely
\(H\), with its scheme structure.

The algebraic fundamental group \(\pi_1(M)\) is
\(\Z\) through \(\nu\): the coroot lattice is the
full cocharacter lattice of a maximal torus of the
simply connected kernel \(H\), and the quotient
cocharacter lattice is that of \(\Gm\). The Hodge
cocharacter has multiplier \(z\) or \(z^{-1}\),
according to the homology convention, so its class
generates this \(\Z\). This proves Hodge maximality
by Cadoret--Moonen's precise criterion
\cite[Sections 2.2--2.3, author p.~4]
{GenCadoretMoonen2020}. Their Theorem A(ii), applied
to the established Mumford--Tate equality of this
finite product, now proves that after a finite
extension \(K\),
\[
 \rho_A(\Gamma_K)\text{ is compact open in }M(\mathbf A_f).
 \tag{E.2.25d}
\]
This is the full restricted-product topology
\cite[Theorem A(ii), author p.~2; Theorem 5.3 and
proof, author pp.~11--12]{GenCadoretMoonen2020}.
The source uses \(H^1\); inverse dual transfers its
actual representation and integral comparison to our
\(H_1\) convention. The multiplier exponent still has
absolute value one, as required above. Further finite
enlargement of \(K\) preserves (E.2.25d), because the
new compact image has finite index in the old one and
is therefore open.

\paragraph{2. Supported conjugation commutators.}
Enlarge the finite set of geometric classes to its
finite \(\Gamma_\Q\)-orbit, choose representatives
\(A_i\), and let actual complex conjugation \(c\)
permute them by \(\pi\). Choose actual isogenies
\(\phi_i:{}^c A_i\to A_{\pi(i)}\). Take \(K/\Q\)
finite Galois defining all models, endomorphisms,
polarizations, these isogenies, and the data of
(E.2.25d). Put
\[
 C_i=(\phi_i)_*c_i:V_i\longrightarrow V_{\pi(i)},
\]
where \(c_i\) is the actual anti-complex-linear
map on homology. Functoriality of Betti--\'etale
comparison gives the exact identity
\[
 C_i\rho_i(c\gamma c)C_i^{-1}
          =\rho_{\pi(i)}(\gamma)
                  \qquad(\gamma\in\Gamma_K).
 \tag{E.2.25e}
\]
These maps identify \(H_i\) with \(H_{\pi(i)}\).
Indeed \(\phi_i\) is a rational Hodge isomorphism,
and \(c_i\) changes the Hodge homomorphism \(h(z)\)
to \(h(\bar z)\); the rational algebraic group
generated by their images is the same. This also
proves the identification when the selected
polarizations are different: their change lies in
the labelled Rosati endomorphism space centralized
by \(H_i\).

Set \(G=\rho(\Gamma_K)\), and let \(\theta\) be
the automorphism induced by \(\gamma\mapsto c\gamma c\)
through (E.2.25e). The compact group
\[
 N=\rho(\Gamma_K\cap\Gamma_{\Q^{\mathrm{cm}}})
\]
is closed and normal in \(G\). It is contained in
\(H(\mathbf A_f)\), since
\(\Q^{\mathrm{ab}}\subset\Q^{\mathrm{cm}}\) and
the multiplier is cyclotomic. Actual \(c\) is
central on \(\Q^{\mathrm{cm}}\). Hence
\[
 [\gamma,c]\in\Gamma_K\cap\Gamma_{\Q^{\mathrm{cm}}},
 \qquad g\theta(g)^{-1}\in N
                \quad(g=\rho(\gamma)\in G).
 \tag{E.2.25f}
\]
By (E.2.25d), \(G\) contains independently supported
open subgroups in every local factor of \(H\), and
the full standard integral factors outside a finite
set of primes. We always use these supported elements.

If \(\theta\) swaps two classes or two local factors,
take \(g\) supported on the first. The element
\(g\theta(g)^{-1}\) has support in the two factors.
Its commutator with \(h\) supported on the first
is \([g,h]\) in that factor and the identity in
the other. Normality of \(N\) supplies this
independently supported derived subgroup.

For a fixed class, we verify nontriviality on every
fixed local simple factor. When \(\End^0(A_i)=\Q\),
the centralizer of the standard symplectic group
is scalar. If \(\operatorname{Ad}(C_i)\) were the
identity, \(C_i\) would be a scalar and thus
complex-linear, contrary to its anti-complex-linearity.
A noncentral rational matrix stays noncentral at
every completion and at all but finitely many
residue characteristics.

For real multiplication let \(\kappa\) be the
automorphism of \(F\) induced by \(C_i\). If it is
nontrivial, it swaps the two \(\operatorname{SL}_2\)
factors at a split prime. At an inert or ramified
prime it acts nontrivially on the centroid \(F_v\)
of the \(\Q_\ell\)-Lie algebra
\(\mathfrak{sl}_2(F_v)\), so the Lie automorphism
is nontrivial. This Lie algebra is \(\Q_\ell\)-simple:
over an algebraic closure its simple summands are
transitively permuted by the local Galois group,
and an ideal is a stable sum of such summands.

If \(\kappa=1\), the matrix \(C_i\) is \(F\)-linear.
On each real \(F\)-component it is an invertible
anti-complex-linear map of a real plane. Such a
matrix has trace zero and square a positive scalar.
Consequently
\(\operatorname{Tr}_F(C_i)=0\) and \(C_i^2=d_i\)
with \(d_i\) totally positive. It is nonscalar
in every \(F_v\), because each \(F\to F_v\) is
injective. Choose a nonzero off-diagonal entry or
diagonal difference in \(F\); excluding the primes
dividing its norm proves nonscalarity in every
residue-field factor outside a finite set.

For quaternionic multiplication, Skolem--Noether
gives \(u\in B^\times\) realizing the automorphism
of the labelled left algebra induced by \(C_i\).
Then \(D_i=u^{-1}C_i\) centralizes \(B\) and is
right multiplication by \(b_i\in(B^{\mathrm{op}})^\times\).
Left multiplication by \(u\) commutes with \(H_i\),
so \(\operatorname{Ad}(C_i)|_{H_i}
=\operatorname{Ad}(D_i)|_{H_i}\). It is a Hodge
endomorphism, so \(D_i\) remains anti-complex-linear.
Under \(B_\R=M_2(\R)\), its right action on the
standard real plane has trace zero and negative
determinant. Thus \(b_i\) is noncentral in every
completion, and in all but finitely many split
matrix reductions. The local Lie algebras of
\(\operatorname{SL}_1(B^{\mathrm{op}})\) are forms
of \(\mathfrak{sl}_2\) and are simple in
characteristic zero.

\paragraph{3. Actual local neighborhoods and the full prime product.}
Consider a fixed local simple factor \(J/\Q_\ell\)
and the supported closed subgroup
\(R=N\cap J(\Q_\ell)\), normal under an open
subgroup \(U\). The preceding argument makes
\(\theta\) nontrivial on \(\Lie J\). For
\(X\in\Lie J\), use the actual Cayley curve
\[
 q_X(t)=(1+tX)(1-tX)^{-1},\qquad q'_X(0)=2X.
 \tag{E.2.25g}
\]
For a symplectic \(X\), the matrices \(1+tX\)
and \(1-tX\) have the same pulled-back alternating
form. For trace-zero \(2\) by \(2\) matrices
their determinants agree, and for a reduced-trace-zero
quaternion their reduced norms agree. Thus (E.2.25g)
is a \(J\)-valued rational analytic curve also at
division primes, and lies in \(U\) on a small ball.

The curve \(q_X(t)\theta(q_X(t))^{-1}\) lies in
\(R\) and has derivative \(2(X-\theta X)\).
The derivatives of its conjugates by \(U\) span
\(\Lie J\). Their span is stable under
\(\operatorname{Ad}U\), hence under the Lie
algebra by differentiating \(\operatorname{Ad}q_Y(t)\);
it is a nonzero ideal in a simple Lie algebra.
Choose finitely many curves with basis derivatives.
Their product has invertible Jacobian at zero.
In local coordinates it is \(J_0x+P(x)\), with
\(P\) of order at least two. On a small ball
\(J_0^{-1}P\) is strictly contracting; iterating
\(x\mapsto y-J_0^{-1}P(x)\) solves
\(J_0x+P(x)=J_0y\) for every sufficiently small
\(y\). Its image contains an actual open
neighborhood of the identity, entirely in \(R\).

The same argument proves that the closed derived
subgroup of an open local semisimple group is open.
Use the curves
\(gq_X(t)g^{-1}q_X(t)^{-1}\), whose derivatives are
\(2(\operatorname{Ad}_g-1)X\). A linear functional
annihilating them annihilates every \([Y,X]\),
by differentiating \(g=q_Y(s)\); semisimplicity
therefore makes their span the entire Lie algebra.
The same product and contraction give an open
neighborhood. For a product, use the supported
curves in its separate simple factors.

Now exclude the finite set containing \(2,3\),
field and quaternion discriminants, denominators
of the fixed maps, noncentrality exceptions, and
the exceptions in (E.2.25d). The supported integral
factors are
\(\operatorname{Sp}_4(\Z_\ell)\),
\(\operatorname{SL}_2(\Z_\ell)\), and
\(\operatorname{SL}_2(\mathcal O_v)\), with
\(\mathcal O_v/\Z_\ell\) unramified quadratic.

For \(\operatorname{Sp}_4\), the reduction of
\(R\) is normal and contains a noncentral
\(g\theta(g)^{-1}\). Otherwise \(\theta\) differs
from the identity by a central character of the
perfect finite symplectic group. That character
is trivial, forcing \(\theta=1\); its standard
centralizer is scalar, contradicting the chosen
noncentral reduction of \(C_i\). The normal
subgroup is therefore the whole finite group
\cite[Lemma 2.3(4), p.~6]{GenMayleWang2024}.
The symplectic lifting lemma then gives
\(R=\operatorname{Sp}_4(\Z_\ell)\)
\cite[Lemma 2.4, p.~6]{GenLombardoSurfaces2015}.

For \(\operatorname{SL}_2(k_v)\), the elementary
simplicity proof in E.2.24 works over every finite
field of characteristic at least five: the faithful
doubly transitive projective-line action, the normal
abelian root group in a point stabilizer, and
perfectness prove simplicity of the projective
quotient. Perfectness rules out a remaining
quotient of order two. A nontrivial semilinear
field automorphism cannot act trivially: traces
of \(\operatorname{SL}_2(k_v)\) cover \(k_v\),
so it would fix every field element; the remaining
matrix centralizer is scalar. Thus the normal
subgroup generated by the \(\theta\)-commutators
has full reduction. Swapped factors have already
been isolated by the supported commutators.

Here is the complete lifting step for the unramified
quadratic ring. Lift each upper unipotent \(U(a)\)
to \(g\in R\), and write
\(g=U(\widetilde a)(1+\ell B)\pmod{\ell^2}\).
The first error in \(g^\ell\) is the sum of
\(\operatorname{Ad}(U(-j\widetilde a))B\) for
\(0\leq j<\ell\). It has degree at most two in
\(j\), and the sums of \(1,j,j^2\) vanish in
\(k_v\). Hence
\[
 g^\ell=1+\ell aE_{12}\pmod{\ell^2}.
\]
All \(a\), the lower roots and their conjugates
span \(\mathfrak{sl}_2(k_v)\) as an
\(\mathbf F_\ell\)-space. The identity
\[
 (1+\ell^nX)^\ell
       =1+\ell^{n+1}X\pmod{\ell^{n+2}}
            \qquad(n\geq1)
\]
supplies every successive congruence layer.
Closedness gives the full integral group.

These integral groups are perfect. For
\(\operatorname{SL}_2\), use
\([
\operatorname{diag}(2,1/2),U(t)]=U(3t)\)
and the lower-root version, with \(3\) a unit.
For \(\operatorname{Sp}_4\), choose a torus
element with root value \(2\) on a short root,
and value \(4\) on a long root. The commutator
coefficient is respectively \(1\) or \(3\).
The root groups generate the integral symplectic
group by elimination: move a unit coordinate of
a primitive first column to its first position,
clear the others by symplectic elementary operations,
clear its paired column, and reduce to
\(\operatorname{Sp}_2\). The same argument proves
the finite-field perfectness used above.

We have therefore obtained independently supported
open subgroups at the exceptional primes and full
integral factors elsewhere. Their finite supported
products lie in \(N\); closedness puts their full
product in \(N\). This proves actual adelic openness.
Finally, for
\(U_0=\Gamma_K\cap\Gamma_{\Q^{\mathrm{cm}}}\),
\[
 \overline{[U_0,U_0]}
       \subset\Gamma_K\cap\Gamma_{F_{\mathrm{cm}}},
 \qquad
 \rho\bigl(\overline{[U_0,U_0]}\bigr)
       =\overline{[N,N]}.
\]
The inclusion holds because
\(L^{\mathrm{ab}}\Q^{\mathrm{cm}}/
\Q^{\mathrm{cm}}\) is abelian for every CM field
\(L\); the equality follows by compactness and
surjectivity onto \(N\). The local open-derived
argument at the exceptional primes and perfectness
elsewhere prove, with the same supported products,
\[
 \rho\bigl(\Gamma_K\cap\Gamma_{F_{\mathrm{cm}}}\bigr)
             \text{ is open in }H(\mathbf A_f).
 \tag{E.2.25h}
\]
Further finite enlargement of \(K\) preserves
this conclusion. Only the displayed derived-subgroup
inclusion was used.

\paragraph{4. The exact rational point and integral coset.}
Every \(H_i\) satisfies the proved criterion D.4.1:
it is simply connected semisimple and has noncompact
real points, respectively \(\operatorname{Sp}_4(\R)\),
\(\operatorname{SL}_2(\R)^2\), or
\(\operatorname{SL}_2(\R)\) in the indefinite
quaternion case. Thus
\[
 H_i(\Q)\text{ is dense in }H_i(\mathbf A_f).
 \tag{E.2.25i}
\]
Fix \(\sigma\in\Gamma_{F_{\mathrm{cm}}}\) and a
conjugate target \({}^{\sigma}A_i\). E.2.11 supplies
an actual rational tensor lift of the prescribed
all-CM and Artin comparison. Its block selected by
the finite-model Artin idempotent is
\[
 T_0:V_i\longrightarrow V_{{}^{\sigma}A_i}.
\]
It intertwines every labelled endomorphism and
polarization and has Tate map \(1\). The actual
completed transport \(\rho_i(\sigma)\) has these
same properties. The centralizers and multiplier-one
kernels in (E.2.25b)--(E.2.25c) therefore give
\[
 a=T_0^{-1}\rho_i(\sigma)\in H_i(\mathbf A_f).
 \tag{E.2.25j}
\]
The underlying group map also has this exact image:
the all-CM kernel of E.2.11 has connected semisimple
image in the finite geometric Mumford--Tate group;
its quotient is a quotient of the geometric CM
protorus. It therefore contains the derived group,
and semisimplicity gives the reverse inclusion.
The image is precisely \(H_i\).

Let \(U\) be an open subgroup of \(H_i(\mathbf A_f)\)
contained in the integral lattice stabilizer.
Choose \(h\in H_i(\Q)\cap aU\) by (E.2.25i).
Then
\[
 T=T_0h\in\rho_i(\sigma)U.
 \tag{E.2.25k}
\]
It is rational, and its completion is an integral
isomorphism between the specified source and target
lattices. Its entries and inverse entries belong to
\(\Q\cap\widehat\Z=\Z\). Thus it is an actual
integral unimodular isomorphism. This is the required
integral intersection of the frozen-target coset;
the rational point used is the established point
\(T_0\) of E.2.11.

\paragraph{5. One simultaneous quartic swap.}
Enumerate one \(A_n\) in every geometric isogeny
class of non-CM simple abelian varieties of dimension
at most two, using the fixed algebraic-model ordering
from E.2.17. Fix bases without an unselected
topological choice as follows. Choose the first
algebraic basis of \(H^0(A_n,\Omega^1)\) in the
same data ordering, also for each conjugate model.
Integration identifies \(\Lambda_{A_n}\) with a
full lattice in \(\C^{\dim A_n}\). Among its
ordered integral bases minimize the sum of the
squared Euclidean norms of their period vectors,
then the lexicographic tuple of all their real
coordinates. Any first available basis supplies
a finite bound for this sum, and only finitely
many ordered lattice bases lie below that bound.
The minimum and its lexicographic tie-break
therefore exist and give a unique basis. No computational
decidability is asserted.

Start with the nonempty compact fibre
\[
 C_0=\{\sigma\in\Gamma_{F_{\mathrm{cm}}}:
                         \sigma|_P=(12)(34)\}
\]
of E.2.6. Choose increasing finite Galois fields
\(K_n\) containing \(P\), the first \(n\) model
data, their finite class orbits and all endomorphisms
and isogenies required in (E.2.25h), and the \(n\)-th
field in the fixed Galois exhaustion of \(\Qb\).
Make each choice by the first suitable item in the
fixed countable data ordering. Partition the
preceding compact fibre by its restriction to
\(K_n\), and choose the first nonempty cylinder
\(C'_n\). All conjugate targets and their bases
are now fixed.

For any \(\sigma_n\in C'_n\), the stabilizer under
right multiplication is exactly
\[
 \mathcal K_n=\Gamma_{F_{\mathrm{cm}}}\cap\Gamma_{K_n}
                   \cap\bigcap_{i<n}\ker\rho_{A_i},
 \qquad C'_n=\sigma_n\mathcal K_n.
 \tag{E.2.25l}
\]
Here \(\mathcal K_n\) is merely notation for this
Galois subgroup. The composition identity
\(\rho_{A_i}(\sigma_n k)
=\rho_{A_i}(\sigma_n)\rho_{A_i}(k)\), valid for
\(k\in\Gamma_{K_n}\), proves the equality of
the stabilizer and the fibre in (E.2.25l).
Write \(J_n\) for the algebraic group \(H_i\)
attached to \(A_n\), to avoid confusing it with
a Galois subgroup. Openness (E.2.25h) for the first
\(n\) factors implies that
\(U_n=\rho_{A_n}(\mathcal K_n)\) is open in
\(J_n(\mathbf A_f)\): its product congruence
neighborhood contains all sufficiently deep
last-factor elements with the preceding components
exactly the identity. The possible next matrices are
\[
 \rho_{A_n}(\sigma_n)U_n.
 \tag{E.2.25m}
\]
By (E.2.25k) this coset contains an actual integral
lattice isomorphism. Choose the first such integer
matrix in entry-height and lexicographic order,
and impose it to obtain \(C_n\subset C'_n\).
These compact nonempty sets are nested. Their
intersection is nonempty and contains only one
\(\sigma_*\), because the prescribed restrictions
to the \(K_n\) exhaust \(\Qb\). This is a
specified Galois element retaining the quartic
action, with integral actual transport on every
\(A_n\).

For this same \(\sigma_*\), every CM factor has
the prescribed integral algebraic comparison of
E.2.6. On a product \(B\) of representatives and
CM factors take the direct sum map. For any
geometric isogeny \(q:B\to A\), put
\[
 L_A=q_*^{-1}\Lambda_A,
 \qquad \Lambda_B\subset L_A\subset\tfrac1m\Lambda_B,
 \qquad L_A/\Lambda_B=\ker(q)\subset B[m](\C)
\]
for a suitable \(m\). Actual completed transport
sends this finite kernel to \(\ker({}^{\sigma_*}q)\).
Its integral source map therefore carries \(L_A\)
onto \(L_{{}^{\sigma_*}A}\), giving the actual
integral map
\[
 T_A=({}^{\sigma_*}q)_*T_Bq_*^{-1}:
       \Lambda_A\xrightarrow{\sim}\Lambda_{{}^{\sigma_*}A}.
\]
Completion injectivity as proved in E.2.16 gives
independence of \(q\) and of the decomposition,
and every correspondence, tensor, exterior-power
and projector equation. Cohomology uses inverse
duals, and the common Tate map is \(1\). The
finite-model Artin blocks have the actual maps
\(e_u\mapsto e_{\sigma_*u}\), so their module
maps give precisely the descent of E.2.3.
This proves (E.2.25a) on the whole specified diagram.

The comparison scope includes every abelian surface
over \(\Qb\). It does not supply integrality for
the additional simple non-CM factors of dimension
at least three or extend the comparison to every
pure motive. E.2.22 remains the exact cover-tower
test for extension of this same \(\sigma_*\) to
the entire abelian diagram.

The exact finite certificate
\path{code/verify_surface_groups.py} checks the
Cayley identities, symplectic root commutators,
the full Lie-bracket span and scalar centralizer,
the real multiplication trace-zero identities,
the quaternion reduced-norm identity, and the
anti-Hodge real-plane matrix. The neighborhood,
prime-product, strong-approximation and simultaneous
comparison assertions are established by the proof
above.
\end{proof}
